\documentclass[11pt,a4paper]{article}

\usepackage[a4paper,margin=25mm]{geometry}
\usepackage{amsmath,amssymb}
\usepackage{siunitx}
\usepackage{booktabs}
\usepackage{circuitikz}
\usepackage{tikz}
\usetikzlibrary{arrows.meta,positioning,calc}
\usepackage[hidelinks]{hyperref}

\ctikzset{bipoles/length=0.9cm}
\sisetup{per-mode=symbol,detect-weight=true}

\newcommand{\dif}{\mathrm{d}}

\title{Energy-Based Control of a Full-Bridge Active Power Decoupling Circuit\\
for a \SI{350}{\kilo\watt} Single-Phase Active Front End}
\author{Design note}
\date{}

\begin{document}
\maketitle

\begin{abstract}
This note develops the control design for a full-bridge active power decoupling
(APD) circuit connected in parallel with the dc link of a single-phase active
front end (AFE). The buffer capacitor is operated with a \emph{bipolar} voltage
at grid frequency, which gives full utilisation of the stored energy and a
bounded, sinusoidal buffer current. The plant is reformulated in terms of stored
energy, which removes the square-root nonlinearity of the voltage reference and
yields two cascaded loops with linear plants. Reference generation, inner
current-loop design with proportional-resonant (PR) compensators, outer energy
loops, stability considerations and a worked design example for
\SI{350}{\kilo\watt} on a \SI{16.7}{\hertz} traction network are given.
\end{abstract}

\section{Scope and problem statement}

A single-phase converter exchanging active power $P$ and reactive power $Q$ with
the grid delivers an instantaneous power
\begin{equation}
p_g(t) = P - S\cos(2\omega_g t + \varphi),
\qquad S=\sqrt{P^2+Q^2},\quad \varphi=\arctan\frac{Q}{P},
\label{eq:pg}
\end{equation}
where $\omega_g$ is the grid angular frequency. The dc load absorbs the constant
term; the oscillating term must be buffered. Integrating it gives the energy
that must be exchanged with the storage element,
\begin{equation}
\Delta E_{pp} = \frac{S}{\omega_g}.
\label{eq:deltaE}
\end{equation}

Equation~\eqref{eq:deltaE} is the governing constraint of the whole design: it
is fixed by the application and cannot be reduced by any choice of topology.
Only the \emph{utilisation} of the storage element and the distribution of the
stress between voltage and current are topology-dependent
\cite{sun2016review,qin2016benchmark}.

For the case considered here ($S\approx P=\SI{350}{\kilo\watt}$,
$f_g=\SI{16.7}{\hertz}$) this gives $\Delta E_{pp}=\SI{3.34}{\kilo\joule}$. A
purely passive dc link limiting the ripple to \SI{5}{\percent} of
$V_{dc}=\SI{700}{\volt}$ would require approximately \SI{136}{\milli\farad},
i.e. of the order of \SI{33}{\kilo\joule} of installed capacitor energy. The
active solution reduces the installed energy to
$\Delta E_{pp}$ itself.

\section{Topology}

The circuit is shown in Fig.~\ref{fig:topology}. A full bridge is connected
across the dc link; its ac terminals feed a series branch made of a small
inductor $L_b$ and a bipolar (film) capacitor bank $C_b$. The bridge is the same
device class as the main converter, since the buffer voltage never exceeds
$V_{dc}$; no pre-charge circuit is required, since the buffer starts from zero
voltage.

\begin{figure}[htb]
\centering
\begin{circuitikz}[scale=0.92,transform shape]
% dc rails
\draw (-4.6,3.2) -- (4.8,3.2);
\draw (-4.6,0.0) -- (4.8,0.0);
% AFE block
\draw[thick] (-6.4,0.4) rectangle (-4.6,2.8);
\node[align=center,font=\footnotesize] at (-5.5,1.6) {1-$\varphi$ AFE\\ + LCL};
\draw (-4.6,2.8) -- (-4.6,3.2);
\draw (-4.6,0.4) -- (-4.6,0.0);
% dc link capacitor
\draw (-3.4,3.2) to[C=$C_{dc}$] (-3.4,0.0);
% load arrow
\draw[-{Latex[length=2mm]}] (-2.0,3.2) -- (-2.0,2.2) node[right,font=\footnotesize] {$i_o$ (load)};
% leg A
\node[nigbt,xscale=-1] (Q1) at (1.0,2.4) {};
\node[nigbt,xscale=-1] (Q2) at (1.0,0.8) {};
\draw (1.0,3.2) -- (Q1.C);
\draw (Q1.E) -- (1.0,1.6);
\draw (1.0,1.6) -- (Q2.C);
\draw (Q2.E) -- (1.0,0.0);
\draw (0.35,1.6) to[D*] (0.35,3.2);
\draw (0.35,0.0) to[D*] (0.35,1.6);
\draw (0.35,1.6) -- (1.0,1.6);
% leg B
\node[nigbt] (Q3) at (4.2,2.4) {};
\node[nigbt] (Q4) at (4.2,0.8) {};
\draw (4.2,3.2) -- (Q3.C);
\draw (Q3.E) -- (4.2,1.6);
\draw (4.2,1.6) -- (Q4.C);
\draw (Q4.E) -- (4.2,0.0);
\draw (4.8,1.6) to[D*] (4.8,3.2);
\draw (4.8,0.0) to[D*] (4.8,1.6);
\draw (4.8,1.6) -- (4.2,1.6);
% buffer branch
\draw (1.0,1.6) to[L=$L_b$,i>^=$i_b$] (2.5,1.6);
\draw (2.5,1.6) to[C=$C_b$,v_=$v_b$] (4.2,1.6);
% labels
\node[font=\footnotesize] at (-2.9,3.5) {$v_{dc}$};
\node[align=center,font=\footnotesize] at (2.6,-0.9) {active power decoupling bridge};
\draw[dashed] (0.0,-0.4) rectangle (5.2,3.7);
\end{circuitikz}
\caption{Full-bridge APD circuit in parallel with the dc link of the
single-phase AFE.}
\label{fig:topology}
\end{figure}

\section{Averaged model and energy formulation}

Let $m\in[-1,1]$ be the averaged modulation index of the bridge. The averaged
model is
\begin{align}
L_b\frac{\dif i_b}{\dif t} &= m\,v_{dc} - v_b - R_b i_b, \label{eq:model_i}\\
C_b\frac{\dif v_b}{\dif t} &= i_b, \label{eq:model_v}\\
C_{dc}\frac{\dif v_{dc}}{\dif t} &= i_{in} - i_o - m\,i_b, \label{eq:model_vdc}
\end{align}
with $R_b$ collecting the branch resistance (capacitor ESR, inductor and
semiconductor conduction losses).

The reference trajectory for $v_b$ is the solution of a quadratic energy balance
and is therefore a square root of a sinusoid; tracking it directly gives a
strongly nonlinear plant \cite{yuan2019nonlinear}. The difficulty disappears
when the buffer energy is taken as the state variable:
\begin{equation}
W_b = \tfrac12 C_b v_b^2 + \tfrac12 L_b i_b^2 .
\end{equation}
Differentiating and substituting \eqref{eq:model_i}--\eqref{eq:model_v},
\begin{equation}
\boxed{\;\frac{\dif W_b}{\dif t} = m\,v_{dc}\,i_b - R_b i_b^2 = p_b - R_b i_b^2\;}
\label{eq:Wb}
\end{equation}
where $p_b$ is the instantaneous power drawn by the APD from the dc link. The
same transformation applied to the dc link, $W_{dc}=\tfrac12 C_{dc}v_{dc}^2$,
gives
\begin{equation}
\frac{\dif W_{dc}}{\dif t} = p_g - p_o - p_b .
\label{eq:Wdc}
\end{equation}
Both \eqref{eq:Wb} and \eqref{eq:Wdc} are integrators driven by power. This is
the basis of the ``power-balance'' family of APD control strategies
\cite{liu2021review,wang2011ripple}.

\section{Steady-state reference trajectories}

Impose a bipolar buffer voltage at \emph{grid} frequency,
\begin{equation}
v_b^\star(t) = \hat V_b\cos\theta, \qquad \theta = \omega_g t + \delta ,
\qquad i_b^\star(t) = -\omega_g C_b \hat V_b \sin\theta .
\label{eq:refs}
\end{equation}
Substituting into $p_b = v_b i_b + L_b i_b \dot i_b$,
\begin{equation}
p_b(t) = -\tfrac12 \omega_g C_b \hat V_b^2 \left(1-\omega_g^2 L_b C_b\right)\sin 2\theta .
\end{equation}
Matching this against the oscillating term of \eqref{eq:pg} gives two conditions:
\begin{equation}
\boxed{\;\delta = \frac{\varphi}{2} + \frac{\pi}{4}\;},
\qquad
\boxed{\;C_b\hat V_b^2\left(1-k_L\right) = \frac{2S}{\omega_g}\;},
\qquad k_L \triangleq \omega_g^2 L_b C_b .
\label{eq:design}
\end{equation}

Three consequences are worth stating explicitly.

\begin{enumerate}
\item \textbf{Phase.} The buffer voltage is a sinusoid at $\omega_g$, phase
shifted by $\varphi/2+\pi/4$ with respect to the grid voltage. Its product with
a current in quadrature produces exactly the required $2\omega_g$ power. No
absolute value, no zero-crossing singularity of the current reference: this is
the structural advantage of the full bridge over the buck-type half bridge
\cite{tang2015symmetrical}.

\item \textbf{Capacitor current.} $\hat I_b = \omega_g C_b \hat V_b =
2S/[\hat V_b(1-k_L)]$. Since $\hat V_b \le V_{dc}$, the current stress is set by
the dc-link voltage. Raising $V_{dc}$ is the single most effective design lever:
$C_b\propto\hat V_b^{-2}$ and $\hat I_b\propto\hat V_b^{-1}$.

\item \textbf{Inductor derating.} The inductor energy oscillates in
anti-phase with the capacitor energy and therefore \emph{subtracts} from the
useful buffering capability by the factor $(1-k_L)$. The branch resonance
$\omega_r = 1/\sqrt{L_bC_b}$ must satisfy $\omega_r \gg \omega_g$; a practical
target is $k_L\le 0.05$, i.e. $\omega_r \ge 4.5\,\omega_g$, while remaining well
below the current-loop bandwidth.
\end{enumerate}

The installed capacitor energy is
\begin{equation}
E_{C_b} = \tfrac12 C_b\hat V_b^2 = \frac{S}{\omega_g\,(1-k_L)} ,
\end{equation}
which for $k_L\to 0$ equals $\Delta E_{pp}$: the utilisation factor is unity,
i.e. the theoretical minimum installed energy is reached.

The required modulation index follows from \eqref{eq:model_i}:
\begin{equation}
\hat m \simeq \frac{\hat V_b\,(1-k_L)}{V_{dc}} \le m_{\max},
\end{equation}
so the inductor voltage partially cancels the capacitor voltage and slightly
relaxes the modulation limit.

\section{Control architecture}

The controller has an inner current loop and two outer energy loops
(Fig.~\ref{fig:control}).

\begin{figure}[htb]
\centering
\resizebox{\textwidth}{!}{%
\begin{tikzpicture}[
  >=Latex,
  block/.style={draw,rectangle,minimum height=9mm,minimum width=16mm,align=center,font=\scriptsize},
  sum/.style={draw,circle,inner sep=1.2pt,font=\tiny},
  lbl/.style={font=\scriptsize}
]
% main chain
\node[sum]   (s1) at (0,0)    {$\Sigma$};
\node[block] (pi) at (2.1,0)  {PI\\(energy)};
\node[sum]   (s2) at (4.2,0)  {$\Sigma$};
\node[block] (rg) at (6.3,0)  {reference\\generator};
\node[sum]   (s3) at (8.4,0)  {$\Sigma$};
\node[block] (pr) at (10.5,0) {PR $\omega_g$\\$(+3\omega_g,5\omega_g)$};
\node[block] (pl) at (12.9,0) {bridge\\$+\;L_b,C_b$};

\draw[->] (s1) -- (pi);
\draw[->] (pi) -- (s2);
\draw[->] (s2) -- (rg);
\draw[->] (rg) -- node[above,lbl]{$i_b^\star$} (s3);
\draw[->] (s3) -- (pr);
\draw[->] (pr) -- (pl);
\draw[->] (pl.east) -- ++(1.0,0) node[right,lbl]{$i_b,v_b$};
\draw[<-] (s1.west) -- ++(-1.2,0) node[left,lbl]{$\overline{W}_b^\star$};

% upper part: feed-forward and 2 omega loop
\node[sum]   (s5)  at (4.2,1.4) {$\Sigma$};
\node[block] (ff)  at (2.0,2.7) {$\hat S,\varphi$ feed-forward\\from AFE};
\node[block] (pr2) at (6.6,2.7) {PR $2\omega_g$};
\node[sum]   (s4)  at (9.2,2.7) {$\Sigma$};

\draw[->] (ff.south) -- (2.0,1.4) -- (s5.west);
\draw[->] (pr2.west) -- (5.3,2.7) -- (5.3,1.4) -- (s5.east);
\draw[->] (s5.south) -- (s2.north);
\draw[->] (s4.west) -- (pr2.east);
\draw[<-] (s4.east) -- ++(1.1,0) node[right,lbl]{$W_{dc}^\star$};
\draw[->] (pl.north) -- (12.9,3.9) -- node[above,lbl]{$W_{dc}$} (9.2,3.9) -- (s4.north);

% lower feedbacks
\draw[->] (12.55,-0.45) -- (12.55,-1.6) -- node[below,lbl]{$i_b$} (8.4,-1.6) -- (s3.south);
\draw[->] (13.25,-0.45) -- (13.25,-2.5) -- node[below,lbl]{$\overline{W}_b$} (0,-2.5) -- (s1.south);
\end{tikzpicture}}
\caption{Control structure. Feed-forward sets the bulk of the ripple-power
command; the $2\omega_g$ resonant loop on the dc-link energy removes the
residual; the slow energy loop replenishes the APD losses.}
\label{fig:control}
\end{figure}

\subsection{Inner current loop}

With $v_b$ used as a feed-forward term, the plant seen by the current controller
is $G_i(s) = 1/(sL_b+R_b)$ preceded by the computation and modulation delay,
approximated by $e^{-1.5 s T_s}$ with $T_s$ the sampling period. Since
$i_b^\star$ is sinusoidal at $\omega_g$, a PR compensator in the stationary
frame gives zero steady-state error:
\begin{equation}
G_c(s) = K_p + \sum_{h\in\{1,3,5\}} \frac{2K_{r,h}\,\omega_c s}{s^2 + 2\omega_c s + (h\omega_g)^2}.
\label{eq:PR}
\end{equation}
The damped form is mandatory here: the traction network frequency is not
perfectly constant, and an undamped resonant term would be detuned. A cut-off
$\omega_c = 2\pi\,(0.5\ldots1)\;\si{\radian\per\second}$ is a reasonable
compromise; if the tolerance band is wide, make $\omega_g$ adaptive from the PLL
\cite{teodorescu2006pr}. The harmonic terms at $3\omega_g$ and $5\omega_g$
compensate dead-time distortion.

Proportional gain from the desired crossover $\omega_{ci}$:
\begin{equation}
K_p = \omega_{ci} L_b, \qquad \omega_{ci} \le \frac{2\pi f_{sw}}{10},
\end{equation}
and $K_{r,1}$ is chosen to place the loop gain at $\omega_g$ some
\SIrange{30}{40}{\deci\bel} above unity; $K_{r,1}\approx(20\ldots50)K_p$ is a
usable starting point, verified on the Bode plot of
$G_c(s)G_i(s)e^{-1.5sT_s}$.

Note that $\omega_r=1/\sqrt{L_bC_b}$ lies inside the current loop; the branch
resonance is damped by the loop provided $\omega_{ci}\ge 3\omega_r$.

\subsection{Buffer energy loop}

The APD is lossy, so the mean value of $p_b$ must equal the APD losses,
otherwise $\overline{W}_b$ drifts. A slow PI acting on
$\overline{W}_b - \overline{W}_b^\star$ corrects $\hat V_b$. Linearising
$\overline{W}_b = \tfrac14 C_b \hat V_b^2$,
\begin{equation}
\frac{\partial \overline{W}_b}{\partial \hat V_b} = \tfrac12 C_b\hat V_b .
\end{equation}
A bandwidth of $f_g/10$ is appropriate; the loop must not react within one grid
period. $\overline{W}_b$ is obtained by low-pass filtering
$\tfrac12 C_b v_b^2 + \tfrac12 L_b i_b^2$, or by a notch at $2\omega_g$.

\subsection{DC-link energy loop}

Feed-forward alone cannot cancel the ripple: errors in $\hat S$, in $C_b$
(tolerance and voltage dependence of film capacitors), and in the dead-time
model leave a residual $2\omega_g$ component. A PR compensator at $2\omega_g$
acting on $W_{dc}$ (not on $v_{dc}$, so that the plant is the linear integrator
of \eqref{eq:Wdc}) produces a correction to the ripple-power command. This is
the ``harmonic suppression'' idea in the classification of
\cite{liu2021review}.

At \SI{16.7}{\hertz} this loop is demanding. The residual dc-link ripple is
\begin{equation}
\Delta V_{dc,pp} \simeq \frac{P_{res}}{\omega_g V_{dc} C_{dc}} ,
\end{equation}
so with $C_{dc}=\SI{5}{\milli\farad}$ and $V_{dc}=\SI{700}{\volt}$ a
\SI{1}{\percent} residual ripple power (\SI{3.5}{\kilo\watt}) already produces
about \SI{9.5}{\volt} peak-to-peak. An attenuation of at least
\SI{40}{\deci\bel} at $2\omega_g$ must therefore be verified on the sensitivity
function, not merely assumed.

\subsection{Reference generation}

\begin{enumerate}
\item Obtain $\hat S$ and $\varphi$ from the AFE current controller (they are
already available as the amplitude and phase of the grid-current reference), and
$\theta_g$ from the PLL.
\item Add the correction from the $2\omega_g$ loop to obtain the effective
ripple-power command $\hat S^\star$.
\item Compute $\hat V_b = \sqrt{2\hat S^\star / [\omega_g C_b(1-k_L)]}$, limited
to $\hat V_{b,\max}$, and $\delta = \varphi/2 + \pi/4$.
\item Generate $v_b^\star$, $i_b^\star$ from \eqref{eq:refs} using
$\theta = \theta_g + \delta$.
\end{enumerate}

Because $\hat V_b \propto \sqrt{\hat S}$, the modulation index falls at partial
load and the relative dead-time distortion grows. Below roughly \SI{30}{\percent}
load the dc-link capacitor alone is usually sufficient, and the APD can be
disabled; the transition must be ramped over several grid periods.

\section{Stability considerations}

\paragraph{Inner loop.} With $\omega_{ci}=2\pi\cdot\SI{300}{\hertz}$ and a
double-update scheme at $f_s=2f_{sw}$, the delay contributes
$1.5\,\omega_{ci}T_s$ of phase lag. The phase margin follows directly from the
Bode plot of $G_c(s)G_i(s)e^{-1.5sT_s}$; the resonant terms reduce the margin
locally and must be included in the plot, not added afterwards.

\paragraph{Cascade separation.} The three loops are separated by more than a
decade in bandwidth ($\approx\SI{300}{\hertz}$, $2\omega_g$, $f_g/10$), so a
sequential design is legitimate. The dc-link \emph{voltage} loop of the main AFE
must either be slower than $2\omega_g/10$ or include a notch at $2\omega_g$;
otherwise it competes with the APD for the same ripple and the two controllers
fight each other.

\paragraph{Periodic nature of the plant.} The modulation index
$m(t)\approx v_b(t)/v_{dc}$ varies sinusoidally at $\omega_g$ with a large
amplitude, so the small-signal model is linear time-periodic. Classical LTI
margins computed at one operating point are indicative only. Two practical
checks: (i) evaluate the loop gain at the extremes of the trajectory
($|m|\to \hat m$ and $m\to 0$) and take the worst case; (ii) verify the Floquet
multipliers of the periodic closed loop, or equivalently run a harmonic
state-space / time-domain sweep. Lyapunov-based designs avoid the issue
altogether by construction \cite{liu2024lyapunov,yuan2019nonlinear}.

\paragraph{DC-link interaction.} Seen from the dc node, the APD behaves as a
controlled power sink with a strong $2\omega_g$ component. An impedance-based
(Nyquist) check of the source--load pair at the dc link, including the AFE
voltage loop and the downstream inverter, is recommended before commissioning.

\paragraph{Saturation.} If $\hat V_b(1-k_L)$ approaches $m_{\max}V_{dc}$ the
bridge loses controllability over part of the period. The reference generator
must limit $\hat V_b$ and the outer loops must have anti-windup; the graceful
degradation is a partial decoupling with the residual ripple absorbed by
$C_{dc}$ \cite{neumayr2021partial}.

\section{Design example}

$S=\SI{350}{\kilo\volt\ampere}$, $f_g=\SI{16.7}{\hertz}$
($\omega_g=\SI{104.9}{\radian\per\second}$), $V_{dc}=\SI{700}{\volt}$,
$\hat V_b=\SI{600}{\volt}$, $L_b=\SI{200}{\micro\henry}$,
$f_{sw}=\SI{3}{\kilo\hertz}$ with unipolar (three-level) PWM.

Solving \eqref{eq:design} for $C_b$ gives the values of
Table~\ref{tab:design}. The last column repeats the exercise for a
\SI{1100}{\volt} dc link, which is still within the same
\SI{1700}{\volt} device class.

\begin{table}[htb]
\centering
\caption{Design example.}
\label{tab:design}
\begin{tabular}{lccc}
\toprule
Quantity & Symbol & $V_{dc}=\SI{700}{\volt}$ & $V_{dc}=\SI{1100}{\volt}$ \\
\midrule
Buffer voltage amplitude & $\hat V_b$ & \SI{600}{\volt} & \SI{950}{\volt} \\
Buffer capacitance & $C_b$ & \SI{19.4}{\milli\farad} & \SI{7.5}{\milli\farad} \\
Inductor factor & $k_L$ & 0.043 & 0.017 \\
Branch resonance & $f_r$ & \SI{81}{\hertz} & \SI{130}{\hertz} \\
Peak buffer current & $\hat I_b$ & \SI{1218}{\ampere} & \SI{750}{\ampere} \\
RMS buffer current & $I_{b,rms}$ & \SI{861}{\ampere} & \SI{530}{\ampere} \\
Installed capacitor energy & $E_{C_b}$ & \SI{3.48}{\kilo\joule} & \SI{3.39}{\kilo\joule} \\
Peak modulation index & $\hat m$ & 0.82 & 0.85 \\
Current ripple, pk-pk & $\Delta I_{pp}$ & \SI{146}{\ampere} & \SI{229}{\ampere} \\
Bridge apparent power & $\hat V_b\hat I_b/2$ & \SI{365}{\kilo\volt\ampere} & \SI{356}{\kilo\volt\ampere} \\
\bottomrule
\end{tabular}

\vspace{2mm}
{\footnotesize Both columns use $L_b=\SI{200}{\micro\henry}$ and
$f_{sw}=\SI{3}{\kilo\hertz}$. For the \SI{1100}{\volt} case the current
ripple is \SI{30}{\percent} of $\hat I_b$; raising $L_b$ to
\SI{400}{\micro\henry} brings it to \SI{115}{\ampere} at the cost of
$k_L=0.034$ and $f_r=\SI{91}{\hertz}$.}
\end{table}

Controller parameters for the \SI{700}{\volt} case:
$\omega_{ci}=2\pi\cdot\SI{300}{\radian\per\second}$,
$K_p=\omega_{ci}L_b=\SI{0.377}{\ohm}$, $K_{r,1}=30K_p$,
$\omega_c=2\pi\cdot\SI{1}{\radian\per\second}$; buffer energy loop at
\SI{1.7}{\hertz}; dc-link resonant loop at $2\omega_g=\SI{33.3}{\hertz}$ with
$\ge\SI{40}{\deci\bel}$ of loop gain at that frequency.

Two observations on the numbers. First, the bridge apparent power is essentially
equal to $S$: an APD in parallel processes the full pulsating power, and its
silicon is comparable to that of the main bridge. Second, the buffer rms current
makes the capacitor ESR a first-order design constraint: at
\SI{861}{\ampere} rms, an equivalent series resistance of \SI{1}{\milli\ohm}
already dissipates about \SI{740}{\watt} in the bank alone.

\section{Practical notes}

\begin{itemize}
\item \textbf{Start-up.} $v_b$ starts from zero; ramp $\hat V_b$ over ten or
more grid periods (\SI{0.6}{\second} at \SI{16.7}{\hertz}). No pre-charge
circuit and no minimum-voltage constraint, unlike a boost-type buffer.

\item \textbf{Dead time.} $v_b$ crosses zero when $i_b$ is at its peak, so the
current polarity is unambiguous where the modulation index is small: the
dead-time error appears as an almost constant voltage offset and is
straightforward to compensate. The residual shows up at $3\omega_g$ and
$5\omega_g$ and is handled by the corresponding resonant terms in
\eqref{eq:PR}.

\item \textbf{Capacitor technology.} The bank must be non-polarised: metallised
polypropylene film. Specify on rms current and $\dif v/\dif t$, and account for
the capacitance drift with temperature and ageing, which the $2\omega_g$ loop
must absorb.

\item \textbf{Interleaving.} At \SI{1218}{\ampere} peak, two interleaved
full-bridge units halve the device current and double the effective ripple
frequency. This is usually cheaper than a single very large module and adds
partial redundancy.

\item \textbf{Fault handling.} On loss of the APD the dc-link ripple returns to
the passive value. The supervisory logic should derate the transferred power
rather than trip, since $\Delta V_{dc}\propto P$.

\item \textbf{Partial decoupling.} Allowing \SIrange{10}{15}{\percent} ripple on
$C_{dc}$ and sizing the APD only for the remainder reduces its rating roughly in
proportion \cite{neumayr2021partial}.
\end{itemize}

\section{Alternative control formulations}

\begin{itemize}
\item \textbf{Virtual impedance / resonator emulation.} The APD is controlled so
that it emulates an LC resonator tuned at $2\omega_g$ seen from the dc link
\cite{liu2021lcresonator}. No explicit estimate of $\hat S$ is needed, which
makes it attractive when the turbine power varies quickly; the trade-off is a
slower transient response.

\item \textbf{Nonlinear and Lyapunov-based control.} Direct treatment of the
nonlinear voltage reference with guaranteed convergence
\cite{yuan2019nonlinear,liu2024lyapunov}; also the natural route to a
passivity-based (IDA-PBC) formulation, for which \eqref{eq:Wb} is already the
required energy function.

\item \textbf{Predictive control.} The energy formulation provides a
one-dimensional cost function ($W_b$ tracking) that suits finite-control-set
MPC, at the cost of a variable switching frequency, which is a drawback for the
EMC qualification of a traction-connected converter.

\item \textbf{Series voltage compensator.} A different architecture altogether:
the compensator is inserted in series with the dc link and only synthesises the
ripple voltage, which lowers its VA rating at the price of carrying the full dc
current \cite{wang2014series}.
\end{itemize}

\begin{thebibliography}{99}

\bibitem{sun2016review}
Y.~Sun, Y.~Liu, M.~Su, W.~Xiong, and J.~Yang, ``Review of active power
decoupling topologies in single-phase systems,'' \emph{IEEE Trans. Power
Electron.}, vol.~31, no.~7, pp.~4778--4794, Jul. 2016.

\bibitem{liu2021review}
Y.~Liu, W.~Zhang, Y.~Sun, M.~Su, G.~Xu, and H.~Dan, ``Review and comparison of
control strategies in active power decoupling,'' \emph{IEEE Trans. Power
Electron.}, vol.~36, no.~12, pp.~14436--14455, Dec. 2021.

\bibitem{wang2011ripple}
R.~Wang, F.~Wang, D.~Boroyevich, R.~Burgos, R.~Lai, P.~Ning, and
K.~Rajashekara, ``A high power density single-phase PWM rectifier with active
ripple energy storage,'' \emph{IEEE Trans. Power Electron.}, vol.~26, no.~5,
pp.~1430--1443, May 2011.

\bibitem{tang2015symmetrical}
Y.~Tang, F.~Blaabjerg, P.~C. Loh, C.~Jin, and P.~Wang, ``Decoupling of
fluctuating power in single-phase systems through a symmetrical half-bridge
circuit,'' \emph{IEEE Trans. Power Electron.}, vol.~30, no.~4, pp.~1855--1865,
Apr. 2015.

\bibitem{tang2015dual}
Y.~Tang, Z.~Qin, F.~Blaabjerg, and P.~C. Loh, ``A dual voltage control strategy
for single-phase PWM converters with power decoupling function,'' \emph{IEEE
Trans. Power Electron.}, vol.~30, no.~12, pp.~7060--7071, Dec. 2015.

\bibitem{yuan2019nonlinear}
H.~Yuan, S.~Li, W.~Qi, S.-C. Tan, and S.~Y. Hui, ``On nonlinear control of
single-phase converters with active power decoupling function,'' \emph{IEEE
Trans. Power Electron.}, vol.~34, no.~6, pp.~5903--5915, Jun. 2019.

\bibitem{liu2021lcresonator}
Y.~Liu, W.~Zhang, J.~Lin, M.~Su, and X.~Liang, ``Active power decoupling control
for single-phase current source rectifier based on emulating LC resonator,''
\emph{IEEE Trans. Ind. Electron.}, vol.~68, no.~6, pp.~5460--5465, Jun. 2021.

\bibitem{qin2016benchmark}
Z.~Qin, Y.~Tang, P.~C. Loh, and F.~Blaabjerg, ``Benchmark of AC and DC active
power decoupling circuits for second-order harmonic mitigation in kW-scale
single-phase inverters,'' \emph{IEEE J. Emerg. Sel. Topics Power Electron.},
vol.~4, no.~1, pp.~15--25, Mar. 2016.

\bibitem{neumayr2021partial}
D.~Neumayr, G.~C. Knabben, E.~Varescon, D.~Bortis, and J.~W. Kolar,
``Comparative evaluation of a full- and partial-power processing active power
buffer for ultracompact single-phase DC/AC converter systems,'' \emph{IEEE J.
Emerg. Sel. Topics Power Electron.}, vol.~9, no.~2, pp.~1994--2013, 2021.

\bibitem{wang2014series}
H.~Wang, H.~S.-H. Chung, and W.~Liu, ``Use of a series voltage compensator for
reduction of the DC-link capacitance in a capacitor-supported system,''
\emph{IEEE Trans. Power Electron.}, vol.~29, no.~3, pp.~1163--1175, Mar. 2014.

\bibitem{liu2024lyapunov}
Y.~Liu \emph{et al.}, ``Reduced sensor Lyapunov-based control for active power
decoupling circuit,'' \emph{Int. J. Circuit Theory Appl.}, 2024.

\bibitem{teodorescu2006pr}
R.~Teodorescu, F.~Blaabjerg, M.~Liserre, and P.~C. Loh,
``Proportional-resonant controllers and filters for grid-connected
voltage-source converters,'' \emph{IEE Proc. Electr. Power Appl.}, vol.~153,
no.~5, pp.~750--762, Sep. 2006.

\end{thebibliography}

\end{document}
